% !TeX root = ActuarialFormSheet_MBFA-es.tex
% !TeX spellcheck = es_ES

\begin{f}[Notaciones de Tabla de Vida]

Edad \(x\), \(y\), \(z\)...    

\(l_x\) es el número de personas vivas, en relación con una cohorte inicial, a la edad \(x\)

\(\omega\) es el límite de edad de las tablas de mortalidad.

\(d_x=l_x-l_{x+1}\) es el número de personas que mueren entre la edad \(x\) y la edad \(x+1\).

\(q_x\) es la probabilidad de muerte entre las edades de \(x\) y la edad \(x+1\).
\[
\,q_x = d_x / l_x 
\]

\(p_x\) es la probabilidad de que el individuo de edad \(x\) sobreviva a la edad \(x+1\).
\[
\,p_x+q_x=1 
\]

Asimismo,
\(\,_nd_x = d_x + d_{x+1} + \cdots + d_{x+n-1} = l_x - l_{x+n}\) muestra el número de personas que mueren entre la edad \(x\) y la edad \(x+n\).

\(\,_nq_x\) es la probabilidad de muerte entre las edades de \(x\) y la edad \(x+n\).

\[
\,_nq_x = {}_nd_x / l_x
\]
\(\,_np_x\) es la probabilidad de que una persona de edad \(x\) sobreviva la edad \(x+n\).
\[
\,_np_x = l_{x+n} / l_x 
\]


\({}_{m|}q_{x}\), la probabilidad de que el individuo de edad \(x\) muera en el \({m+1}^{th}\) año.
\[{}_{m|}q_{x}=\frac{d_{x+m}}{l_x}=\frac{l_{x+m}-l_{x+m+1}}{l_x}\]

\(\,e_x\) es la esperanza de vida de una persona aún viva a la edad \(x\).
Este es el número de cumpleaños que esperas vivir.
\[
\,e_x = \sum_{t=1}^{\infty} \ _tp_x 
\]
\end{f}
\hrule

\begin{f}[Coeficiente de conmutaciones]


Estos coeficientes o conmutaciones se establecen mediante funciones actuariales que dependen de una tabla de mortalidad y de una tasa \(i\) (\(v=1/(1+i)\)) para establecer la tabla actuarial.
\[
D_x=l_x .v^x
\]
puede verse "como" el número real de supervivientes. Las sumas

\[
N_x=\sum_{k\geq 0} D_{x+k}=\sum_{k= 0}^{\omega-x} D_{x+k}
\]

\[
S_x=\sum_{k\geq 0} N_{x+k}=\sum_{k\geq 0}(k+1). D_{x+k}
\]
se utilizará para simplificar los cálculos.
Asimismo
\[
C_x = d_x v^{ x+1} 
\]
puede verse "como" el número de muertes descontado a la edad \(x\). Las sumas

\[
M_x=\sum_{k= 0}^{\omega-x} C_{x+k}
\]
\[
R_x=\sum_{k= 0}^{\omega-x} M_{x+k}
\]
se utilizará para simplificar los cálculos.

Los coeficientes \(D_x\) \(N_x\) y \(S_x\) se utilizarán para los cálculos de las operaciones en caso de vida y \(C_x\) \(M_x\) y \(R_x\) para las operaciones en caso de muerte.

\end{f} 
\hrule

\begin{f}[Rentas vitalicias o anualidades]


\medskip
	

\begin{tikzpicture}[scale=0.75]
    % Draw the x-axis and y-axis.
    \def\w{11}
    \def\n{6}
    \node[left] at (-.5,0) {\({}_{}a_x\)};
    
    \begin{scope}[shift={(3.75,.25)}]
        \draw[color=OrangeProfondIRA,scale=0.2,fill=OrangeProfondIRA] \Cerceuil;
    \end{scope}
    \foreach \y in  {0,...,3} {
        \draw (\y,0) -- (\y,-0.1);
        \ifthenelse{\y>0 }{	\node[below] at (\y,-0.1) {\tiny \( \scriptstyle x+\y\)};
            \draw[ line width=1, color=OrangeProfondIRA, arrows={-Stealth[length=4, inset=0]}] (\y,0) -- (\y,1);}{
            \node[below] at (\y,-0.1) {\tiny \( \scriptstyle x\)};}
    }
    \draw (\n,0) -- (\n,-0.1);
    \node[below] at (\n,-0.1) {\tiny \(\scriptstyle  x+n\)};
    \foreach \y in  {1,...,3} {
        \draw (\y+\n,0) -- (\y+\n,-0.1);
        \node[below] at (\y+\n,-0.1) {\tiny \(\scriptstyle x+n+\y\)};
        %		\draw[ line width=1, color=OrangeProfondIRA, arrows={-Stealth[length=4, inset=0]}] (\y+\n,0) -- (\y+\n,1);
    }
    \draw[arrows={-Stealth[length=4, inset=0]}, line width=1] (-.5,0) -- (\w,0);
\end{tikzpicture}

\begin{tikzpicture}[scale=0.75]
    % Draw the x-axis and y-axis.
    \def\w{11}
    \def\n{6}
    \node[left] at (-.5,0) {\({}_{}a_x\)};
    
    \begin{scope}[shift={(\n+.5+3,.25)}]
        \draw[color=OrangeProfondIRA,scale=0.2,fill=OrangeProfondIRA] \Cerceuil;
    \end{scope}
    \foreach \y in  {0,...,3} {
        \draw (\y,0) -- (\y,-0.1);
        \ifthenelse{\y>0 }{	\node[below] at (\y,-0.1) {\tiny \( \scriptstyle x+\y\)};
            \draw[ line width=1, color=OrangeProfondIRA, arrows={-Stealth[length=4, inset=0]}] (\y,0) -- (\y,1);}{
            \node[below] at (\y,-0.1) {\tiny \( \scriptstyle x\)};}
    }
    \foreach \y in  {0,...,4} {
        \draw (\y+\n,0) -- (\y+\n,-0.1);
        \ifthenelse{\y>0 }{\node[below] at (\y+\n,-0.1) {\tiny \(\scriptstyle x+n+\y\)};}{
            \node[below] at (\y+\n,-0.1) {\tiny \(\scriptstyle x+n\)};}
        \ifthenelse{\y<4 }{	\draw[ line width=1, color=OrangeProfondIRA, arrows={-Stealth[length=4, inset=0]}] (\y+\n,0) -- (\y+\n,1);}
    }
    \draw[arrows={-Stealth[length=4, inset=0]}, line width=1] (-.5,0) -- (\w,0);
\end{tikzpicture}
\begin{tikzpicture}[scale=0.75]
    % Draw the x-axis and y-axis.
    \def\w{11}
    \def\n{6}
    \node[left] at (-.5,0) {\(\ddot{a}_x\)};
    
    \begin{scope}[shift={(3.75,.25)}]
        \draw[color=OrangeProfondIRA,scale=0.2,fill=OrangeProfondIRA] \Cerceuil;
    \end{scope}
    \foreach \y in  {0,...,3} {
        \draw (\y,0) -- (\y,-0.1);
        \ifthenelse{\y>0 }{	\node[below] at (\y,-0.1) {\tiny \( \scriptstyle x+\y\)};}{
            \node[below] at (\y,-0.1) {\tiny \( \scriptstyle x\)};}
        \draw[ line width=1, color=OrangeProfondIRA, arrows={-Stealth[length=4, inset=0]}] (\y,0) -- (\y,1);
    }
    \draw (\n,0) -- (\n,-0.1);
    \node[below] at (\n,-0.1) {\tiny \(\scriptstyle  x+n\)};
    \foreach \y in  {1,...,4} {
        \draw (\y+\n,0) -- (\y+\n,-0.1);
        \node[below] at (\y+\n,-0.1) {\tiny \(\scriptstyle x+n+\y\)};
        %		\draw[ line width=1, color=OrangeProfondIRA, arrows={-Stealth[length=4, inset=0]}] (\y+\n,0) -- (\y+\n,1);
    }
    \draw[arrows={-Stealth[length=4, inset=0]}, line width=1] (-.5,0) -- (\w,0);
\end{tikzpicture}
\begin{tikzpicture}[scale=0.75]
    % Draw the x-axis and y-axis.
    \def\w{11}
    \def\n{6}
    \node[left] at (-.5,0) {\(\ddot{a}_x\)};
    
    \begin{scope}[shift={(\n+.5+3,.25)}]
        \draw[color=OrangeProfondIRA,scale=0.2,fill=OrangeProfondIRA] \Cerceuil;
    \end{scope}
    \foreach \y in  {0,...,3} {
        \draw (\y,0) -- (\y,-0.1);
        \ifthenelse{\y>0 }{	\node[below] at (\y,-0.1) {\tiny \( \scriptstyle x+\y\)};}{
            \node[below] at (\y,-0.1) {\tiny \( \scriptstyle x\)};}
        \draw[ line width=1, color=OrangeProfondIRA, arrows={-Stealth[length=4, inset=0]}] (\y,0) -- (\y,1);
    }
    \foreach \y in  {0,...,4} {
        \draw (\y+\n,0) -- (\y+\n,-0.1);
        \ifthenelse{\y>0 }{\node[below] at (\y+\n,-0.1) {\tiny \(\scriptstyle x+n+\y\)};}{
            \node[below] at (\y+\n,-0.1) {\tiny \(\scriptstyle x+n\)};}
        \ifthenelse{\y<4 }{	\draw[ line width=1, color=OrangeProfondIRA, arrows={-Stealth[length=4, inset=0]}] (\y+\n,0) -- (\y+\n,1);}
    }
    \draw[arrows={-Stealth[length=4, inset=0]}, line width=1] (-.5,0) -- (\w,0);
\end{tikzpicture}	
\[a_x %=\frac{N_{x+1}}{D_x} 
%=\sum_{k=1}^{\infty}{}_{k|}q_{x} \ddot{a}_{\lcroof{k+1}}
=\sum_{k=1}^{\infty}{}_{k}p_{x} v^{k}=\ddot{a}_x -1
=\frac{N_{x+1}}{D_{x}}
\]

\[\ddot{a}_x 
%\frac{N_x}{D_x}=
%=\sum_{k=0}^{\infty}{}_{k|}q_{x} \ddot{a}_{\lcroof{k+1}}
=\sum_{k=0}^{\infty}{}_{k}p_{x} v^{k}
=	\frac{N_{x}}{D_{x}} 
\]

Si la periodicidad corresponde a \(m\) períodos por año:
\[\ddot{a}_{x}^{(m)} 
=\sum_{k=0}^{\infty}\frac{1}{m}{}_{\frac{k}{m}}p_{x} v^{\frac{k}{m}}\approx\ddot{a}_x -\frac{m-1}{2m}
\]
De manera similar, si paga \(1/m\) al inicio de los \(m\) períodos
\[a_{x}^{(m)}\approx a_x +\frac{m-1}{2m}
\]


\textbf{Rentas vitalicias temporales}. (\engl{Whole life annuity guaranteed for n years})
\[
a_{x:\lcroof{n}} =
\sum_{k=1}^{n}{}_{k}p_{x} v^{k}
=\frac{N_{x+1}-N_{x+n+1}}{D_{x}}
\]
	
\[
\ddot{a}_{x:\lcroof{n}} =%\frac{N_x - N_{x+n}}{D_x}=
\sum_{k=0}^{n-1}{}_{k}p_{x} v^{k}
=\frac{N_{x}-N_{x+n}}{D_{x}}
\]


\textbf{Rentas vitalicias diferidas}
\({}_{m|}a_{x}\) (\engl{Deferred life annuity}) representan las anualidades del individuo de edad \(x\) diferidas \(m\) años. El primer pago se realiza en \(m+1\) años en el caso de la vida.

	%	\includegraphics[width=1\linewidth]{../../LifeActuarial/Graph/RenteViagereDifferee}
\begin{tikzpicture}[scale=0.75]
    % Draw the x-axis and y-axis.
    \def\w{11}
    \def\m{6}
    \node[left] at (-.5,0) {\({}_{m|}a_x\)};
    
    \begin{scope}[shift={(4.25,.25)}]
        \draw[color=OrangeProfondIRA,scale=0.2,fill=OrangeProfondIRA] \Cerceuil;
    \end{scope}
    \draw[dashed, color=BleuProfondIRA,arrows={Stealth[length=4, inset=0]-Stealth[length=4, inset=0]},  line width=1] (0,.3) -- (\m,.3) node [pos=0.5, above] {\(m\)};		
    \draw (0,0) -- (0,-0.1);
    \node[below] at (0,-0.1) {\tiny \(x\)};
    \foreach \y in  {1,...,3} {
        \draw (\y,0) -- (\y,-0.1);
        \node[below] at (\y,-0.1) {\tiny \( \scriptstyle x+\y\)};
    }
    \draw (\m,0) -- (\m,-0.1);
    \node[below] at (\m,-0.1) {\tiny \(\scriptstyle  x+m\)};
    \foreach \y in  {1,...,3} {
        \draw (\y+\m,0) -- (\y+\m,-0.1);
        \node[below] at (\y+\m,-0.1) {\tiny \(\scriptstyle x+m+\y\)};
        %		\draw[ line width=1, color=OrangeProfondIRA, arrows={-Stealth[length=4, inset=0]}] (\y+\m,0) -- (\y+\m,1);
        \draw[arrows={-Stealth[length=4, inset=0]}, line width=1] (-.5,0) -- (\w,0);
    }
\end{tikzpicture}

\begin{tikzpicture}[scale=0.75]
    % Draw the x-axis and y-axis.
    \def\w{11}
    \def\m{6}
    \node[left] at (-.5,0) {\({}_{m|}a_x\)};
    
    \begin{scope}[shift={(\m+.5+3,.25)}]
        \draw[color=OrangeProfondIRA,scale=0.2,fill=OrangeProfondIRA] \Cerceuil;
    \end{scope}
    \draw[dashed, color=BleuProfondIRA,arrows={Stealth[length=4, inset=0]-Stealth[length=4, inset=0]},  line width=1] (0,.3) -- (\m,.3) node [pos=0.5, above] {\(m\)};		
    \draw (0,0) -- (0,-0.1);
    \node[below] at (0,-0.1) {\tiny \(x\)};
    \foreach \y in  {1,...,3} {
        \draw (\y,0) -- (\y,-0.1);
        \node[below] at (\y,-0.1) {\tiny \( \scriptstyle x+\y\)};
    }
    \draw (\m,0) -- (\m,-0.1);
    \node[below] at (\m,-0.1) {\tiny \(\scriptstyle  x+m\)};
    \foreach \y in  {1,...,3} {
        \draw (\y+\m,0) -- (\y+\m,-0.1);
        \node[below] at (\y+\m,-0.1) {\tiny \(\scriptstyle x+m+\y\)};
        \draw[ line width=1, color=OrangeProfondIRA, arrows={-Stealth[length=4, inset=0]}] (\y+\m,0) -- (\y+\m,1);
    }
    \draw[arrows={-Stealth[length=4, inset=0]}, line width=1] (-.5,0) -- (\w,0);
\end{tikzpicture}


  
\end{f}
\hrule

\begin{f}[Prestaciones por muerte o supervivencia]

% !TeX spellcheck = es_ES
\textbf{Beneficios por muerte}(\engl{Whole life insurance } anotado \({SP}_{x}\) or \({A}_{x}\))

\(A_x\) indica un beneficio por muerte al final del año de la muerte (monto de 1), independientemente de la fecha de ocurrencia, para una persona asegurada a la edad \(x\) en el momento de la suscripción.

\(A_{x:\lcroof{n}}\) denota un capital pagado en caso de fallecimiento si éste ocurre y a más tardar en \(n\) años (\engl{Endowment}).

\(\lcterm{A}{x}{n}\) % ou \(\termins{x}{n}\) 
denota un beneficio por muerte pagado si \(x\) muere dentro de los próximos \(n\) años (\engl{Term insurance}).


\(A_x^{(12)}\) indica un beneficio pagadero al final del mes de la muerte.

\(\overline{A}_x\) Indica un beneficio pagado en la fecha de muerte.
\begin{tikzpicture}[scale=0.75]
% Draw the x-axis and y-axis.
\def\w{11}
\def\n{7}
\node[left] at (-.5,0) {\(A_x\)};

\begin{scope}[shift={(3.75,.25)}]
    \draw[color=OrangeProfondIRA,scale=0.2,fill=OrangeProfondIRA] \Cerceuil;
\end{scope}
\draw[ line width=1, color=OrangeProfondIRA, arrows={-Stealth[length=4, inset=0]}] (4,0) -- (4,1);
\foreach \y in  {0,...,4} {
    \draw (\y,0) -- (\y,-0.1);
    \ifthenelse{\y>0 }{	\node[below] at (\y,-0.1) {\tiny \( \scriptstyle x+\y\)};}{
        \node[below] at (\y,-0.1) {\tiny \( \scriptstyle x\)};}
}
\draw (\n,0) -- (\n,-0.1);
\node[below] at (\n,-0.1) {\tiny \(\scriptstyle  x+n\)};
\foreach \y in  {1,...,3} {
    \draw (\y+\n,0) -- (\y+\n,-0.1);
    \node[below] at (\y+\n,-0.1) {\tiny \(\scriptstyle x+n+\y\)};
    %		\draw[ line width=1, color=OrangeProfondIRA, arrows={-Stealth[length=4, inset=0]}] (\y+\n,0) -- (\y+\n,1);
}
\draw[arrows={-Stealth[length=4, inset=0]}, line width=1] (-.5,0) -- (\w,0);
\end{tikzpicture}
\begin{tikzpicture}[scale=0.75]
% Draw the x-axis and y-axis.
\def\w{11}
\def\n{7}
\node[left] at (-.5,0) {\(A_x\)};

\begin{scope}[shift={(\n+.5+2,.25)}]
    \draw[color=OrangeProfondIRA,scale=0.2,fill=OrangeProfondIRA] \Cerceuil;
\end{scope}
\draw[ line width=1, color=OrangeProfondIRA, arrows={-Stealth[length=4, inset=0]}] (\n+3,0) -- (\n+3,1);
\foreach \y in  {0,...,4} {
    \draw (\y,0) -- (\y,-0.1);
    \ifthenelse{\y>0 }{	\node[below] at (\y,-0.1) {\tiny \( \scriptstyle x+\y\)};}{
        \node[below] at (\y,-0.1) {\tiny \( \scriptstyle x\)};}
}
\draw (\n,0) -- (\n,-0.1);
\node[below] at (\n,-0.1) {\tiny \(\scriptstyle  x+n\)};
\foreach \y in  {1,...,3} {
    \draw (\y+\n,0) -- (\y+\n,-0.1);
    \node[below] at (\y+\n,-0.1) {\tiny \(\scriptstyle x+n+\y\)};
    %		\draw[ line width=1, color=OrangeProfondIRA, arrows={-Stealth[length=4, inset=0]}] (\y+\n,0) -- (\y+\n,1);
}
\draw[arrows={-Stealth[length=4, inset=0]}, line width=1] (-.5,0) -- (\w,0);
\end{tikzpicture}
\begin{tikzpicture}[scale=0.75]
% Draw the x-axis and y-axis.
\def\w{11}
\def\n{7}
\node[left] at (-.5,0) {\(\overline{A}_x\)};

\begin{scope}[shift={(3.75,.25)}]
    \draw[color=OrangeProfondIRA,scale=0.2,fill=OrangeProfondIRA] \Cerceuil;
\end{scope}
\draw[ line width=1, color=black, arrows={-Stealth[length=4, inset=0]}] (3.75,0) -- (3.75,1);
\foreach \y in  {0,...,4} {
    \draw (\y,0) -- (\y,-0.1);
    \ifthenelse{\y>0 }{	\node[below] at (\y,-0.1) {\tiny \( \scriptstyle x+\y\)};}{
        \node[below] at (\y,-0.1) {\tiny \( \scriptstyle x\)};}
}
\draw (\n,0) -- (\n,-0.1);
\node[below] at (\n,-0.1) {\tiny \(\scriptstyle  x+n\)};
\foreach \y in  {1,...,3} {
    \draw (\y+\n,0) -- (\y+\n,-0.1);
    \node[below] at (\y+\n,-0.1) {\tiny \(\scriptstyle x+n+\y\)};
    %		\draw[ line width=1, color=OrangeProfondIRA, arrows={-Stealth[length=4, inset=0]}] (\y+\n,0) -- (\y+\n,1);
}
\draw[arrows={-Stealth[length=4, inset=0]}, line width=1] (-.5,0) -- (\w,0);
\end{tikzpicture}
\begin{tikzpicture}[scale=0.75]
% Draw the x-axis and y-axis.
\def\w{11}
\def\n{7}
\node[left] at (-.5,0) {\(\overline{A}_x\)};

\begin{scope}[shift={(\n+.5+2,.25)}]
    \draw[color=OrangeProfondIRA,scale=0.2,fill=OrangeProfondIRA] \Cerceuil;
\end{scope}
\draw[ line width=1, color=black, arrows={-Stealth[length=4, inset=0]}] (\n+.5+2,0) -- (\n+.5+2,1);
\foreach \y in  {0,...,4} {
    \draw (\y,0) -- (\y,-0.1);
    \ifthenelse{\y>0 }{	\node[below] at (\y,-0.1) {\tiny \( \scriptstyle x+\y\)};}{
        \node[below] at (\y,-0.1) {\tiny \( \scriptstyle x\)};}
}
\draw (\n,0) -- (\n,-0.1);
\node[below] at (\n,-0.1) {\tiny \(\scriptstyle  x+n\)};
\foreach \y in  {1,...,3} {
    \draw (\y+\n,0) -- (\y+\n,-0.1);
    \node[below] at (\y+\n,-0.1) {\tiny \(\scriptstyle x+n+\y\)};
    %		\draw[ line width=1, color=OrangeProfondIRA, arrows={-Stealth[length=4, inset=0]}] (\y+\n,0) -- (\y+\n,1);
}
\draw[arrows={-Stealth[length=4, inset=0]}, line width=1] (-.5,0) -- (\w,0);
\end{tikzpicture}

\medskip

Beneficio de vida entera (\engl{Whole life})
\[A_{x}=\sum_{k=0}^{\infty} {}_{k|}q_x\ \nu^{k+1}=\frac{M_x}{D_x}\]
	
\[\termins{x}{n}=\sum_{k=0}^{n-1} {}_{k|}q_x\ \nu^{k+1}=\frac{M_x-M_{x+n}}{D_x}
\]



\medskip    
\textbf{Capital diferido (\engl{Pure Endowment}, capital único en caso de supervivencia)} señaló \(\lcend{A}{x}{n}\) %\(\pureend{x}{n}\)
o \({}_n E_x\).

\begin{tikzpicture}[scale=0.75]
% Draw the x-axis and y-axis.
\def\w{11}
\def\n{7}
\node[left] at (-.5,0) {\({}_n E_x\)};

\begin{scope}[shift={(3.75,.25)}]
    \draw[color=OrangeProfondIRA,scale=0.2,fill=OrangeProfondIRA] \Cerceuil;
\end{scope}
%		\draw[ line width=1, color=OrangeProfondIRA, arrows={-Stealth[length=4, inset=0]}] (4,0) -- (4,1);
\foreach \y in  {0,...,4} {
    \draw (\y,0) -- (\y,-0.1);
    \ifthenelse{\y>0 }{	\node[below] at (\y,-0.1) {\tiny \( \scriptstyle x+\y\)};}{
        \node[below] at (\y,-0.1) {\tiny \( \scriptstyle x\)};}
}
\draw (\n,0) -- (\n,-0.1);
\node[below] at (\n,-0.1) {\tiny \(\scriptstyle  x+n\)};
\foreach \y in  {1,...,3} {
    \draw (\y+\n,0) -- (\y+\n,-0.1);
    \node[below] at (\y+\n,-0.1) {\tiny \(\scriptstyle x+n+\y\)};
    %		\draw[ line width=1, color=OrangeProfondIRA, arrows={-Stealth[length=4, inset=0]}] (\y+\n,0) -- (\y+\n,1);
}
\draw[arrows={-Stealth[length=4, inset=0]}, line width=1] (-.5,0) -- (\w,0);
\end{tikzpicture}
\begin{tikzpicture}[scale=0.75]
% Draw the x-axis and y-axis.
\def\w{11}
\def\n{7}
\node[left] at (-.5,0) {\({}_n E_x\)};

\begin{scope}[shift={(\n+.5+2,.25)}]
    \draw[color=OrangeProfondIRA,scale=0.2,fill=OrangeProfondIRA] \Cerceuil;
\end{scope}
\draw[ line width=1, color=OrangeProfondIRA, arrows={-Stealth[length=4, inset=0]}] (\n,0) -- (\n,1);
\foreach \y in  {0,...,4} {
    \draw (\y,0) -- (\y,-0.1);
    \ifthenelse{\y>0 }{	\node[below] at (\y,-0.1) {\tiny \( \scriptstyle x+\y\)};}{
        \node[below] at (\y,-0.1) {\tiny \( \scriptstyle x\)};}
}
\draw (\n,0) -- (\n,-0.1);
\node[below] at (\n,-0.1) {\tiny \(\scriptstyle  x+n\)};
\foreach \y in  {1,...,3} {
    \draw (\y+\n,0) -- (\y+\n,-0.1);
    \node[below] at (\y+\n,-0.1) {\tiny \(\scriptstyle x+n+\y\)};
    %		\draw[ line width=1, color=OrangeProfondIRA, arrows={-Stealth[length=4, inset=0]}] (\y+\n,0) -- (\y+\n,1);
}
\draw[arrows={-Stealth[length=4, inset=0]}, line width=1] (-.5,0) -- (\w,0);
\end{tikzpicture}
	\[
	={}_n E_x={}_n p_x .v^n=\frac{l_{x+n}}{l_x} . v^n =\frac{D_{x+n}}{D_{x}}
	\]
Prestación por fallecimiento con pago del capital en caso de supervivencia (\engl{Endowment})
\[A_{x:\lcroof{n}}=\termins{x}{n}+\lcend{A}{x}{n}\]

\end{f}
\hrule

\begin{f}[Seguro de vida para varias personas] 

\(a_{xyz}\) es una anualidad, pagada al final del primer año y durante el tiempo que vivan \((x)\), \((y)\) y \((z)\).

\(a_{\overline{xyz}}\) es una anualidad, pagada al final del primer año y durante el tiempo que vivan \((x)\), \((y)\) o \((z)\).

\[
a_{\overline{xy}}=a_{y}+a_{x}-a_{xy}
\]

\(A_{xyz}\) es un seguro que entra en vigor al final del año del primer fallecimiento de \((x)\), \((y)\) y \((z)\).

La barra vertical indica condicionalidad :

\(a_{x|y}\) es una anualidad de sobreviviente que beneficia a \((x)\) después de la muerte de \((y)\).

\(A_{x|yz}\) es un seguro de primero en morir \((y)\) y \((z)\).	


\begin{tikzpicture}[scale=0.85]
    % Draw the x-axis and y-axis.
    \def\w{11}
    \def\n{6}
    \node[left] at (-.5,0) {\(a_{{\color{OrangeProfondIRA}x}|{\color{BleuProfondIRA}y}}\)};
    
    \begin{scope}[shift={(9.75,.25)}]
        \draw[color=OrangeProfondIRA,scale=0.2,fill=OrangeProfondIRA] \Cerceuil;
    \end{scope}
    \begin{scope}[shift={(7.55,.25)}]
        \draw[color=BleuProfondIRA,scale=0.2,fill=BleuProfondIRA] \Cerceuil;
    \end{scope}
    \draw[ line width=1, color=OrangeProfondIRA, arrows={-Stealth[length=4, inset=0]}] (8,0) -- (8,1);
    \draw[ line width=1, color=OrangeProfondIRA, arrows={-Stealth[length=4, inset=0]}] (9,0) -- (9,1);
    \foreach \y in  {0,...,3} {
        \draw (\y,0) -- (\y,-0.1);
        \ifthenelse{\y>0 }{	\node[below] at (\y,-0.1) {\tiny \( \scriptstyle x+\y\)};}{
            \node[below] at (\y,-0.1) {\tiny \( \scriptstyle x\)};}
    }
    \draw (\n,0) -- (\n,-0.1);
    \node[below] at (\n,-0.1) {\tiny \(\scriptstyle  x+n\)};
    \foreach \y in  {1,...,3} {
        \draw (\y+\n,0) -- (\y+\n,-0.1);
        \node[below] at (\y+\n,-0.1) {\tiny \(\scriptstyle x+n+\y\)};
    }
    \draw[arrows={-Stealth[length=4, inset=0]}, line width=1] (-.5,0) -- (\w,0);
\end{tikzpicture}
\begin{tikzpicture}[scale=0.85]
    % Draw the x-axis and y-axis.
    \def\w{11}
    \def\n{6}
    \node[left] at (-.5,0) {\(a_{{\color{OrangeProfondIRA}x}|{\color{BleuProfondIRA}y}}\)};
    
    \begin{scope}[shift={(8.4,.25)}]
        \draw[color=OrangeProfondIRA,scale=0.2,fill=OrangeProfondIRA] \Cerceuil;
    \end{scope}
    \begin{scope}[shift={(2.75,.25)}]
        \draw[color=BleuProfondIRA,scale=0.2,fill=BleuProfondIRA] \Cerceuil;
    \end{scope}
    \foreach \y in  {3,...,8} {
        \draw[ line width=1, color=OrangeProfondIRA, arrows={-Stealth[length=4, inset=0]}] (\y,0) -- (\y,1);
    }
    \foreach \y in  {0,...,3} {
        \draw (\y,0) -- (\y,-0.1);
        \ifthenelse{\y>0 }{	\node[below] at (\y,-0.1) {\tiny \( \scriptstyle x+\y\)};}{
            \node[below] at (\y,-0.1) {\tiny \( \scriptstyle x\)};}
    }
    \foreach \y in  {0,...,4} {
        \draw (\y+\n,0) -- (\y+\n,-0.1);
        \ifthenelse{\y>0 }{\node[below] at (\y+\n,-0.1) {\tiny \(\scriptstyle x+n+\y\)};}{
            \node[below] at (\y+\n,-0.1) {\tiny \(\scriptstyle x+n\)};}
    }
    \draw[arrows={-Stealth[length=4, inset=0]}, line width=1] (-.5,0) -- (\w,0);
\end{tikzpicture}
\begin{tikzpicture}[scale=0.85]
    % Draw the x-axis and y-axis.
    \def\w{11}
    \def\n{6}
    \node[left] at (-.5,0) {\(a_{{\color{OrangeProfondIRA}x}|{\color{BleuProfondIRA}y}}\)};
    
    \begin{scope}[shift={(3.75,.25)}]
        \draw[color=OrangeProfondIRA,scale=0.2,fill=OrangeProfondIRA] \Cerceuil;
    \end{scope}
    \begin{scope}[shift={(7.55,.25)}]
        \draw[color=BleuProfondIRA,scale=0.2,fill=BleuProfondIRA] \Cerceuil;
    \end{scope}
    \foreach \y in  {0,...,3} {
        \draw (\y,0) -- (\y,-0.1);
        \ifthenelse{\y>0 }{	\node[below] at (\y,-0.1) {\tiny \( \scriptstyle x+\y\)};}{
            \node[below] at (\y,-0.1) {\tiny \( \scriptstyle x\)};}
    }
    \draw (\n,0) -- (\n,-0.1);
    \node[below] at (\n,-0.1) {\tiny \(\scriptstyle  x+n\)};
    \foreach \y in  {1,...,4} {
        \draw (\y+\n,0) -- (\y+\n,-0.1);
        \node[below] at (\y+\n,-0.1) {\tiny \(\scriptstyle x+n+\y\)};
        %		\draw[ line width=1, color=OrangeProfondIRA, arrows={-Stealth[length=4, inset=0]}] (\y+\n,0) -- (\y+\n,1);
    }
    \draw[arrows={-Stealth[length=4, inset=0]}, line width=1] (-.5,0) -- (\w,0);
\end{tikzpicture}
\end{f}

\begin{f}[Esquemas simplificados de precios para primas periódicas y reservas]
	
\tikzstyle{startstop} = [rectangle, rounded corners, text width=4cm, minimum height=1cm,text centered, draw=black, fill=BleuProfondIRA!40]
\tikzstyle{process} = [rectangle, minimum width=4cm, minimum height=1cm, text centered, text width=3cm, draw=black, fill=OrangeProfondIRA!40]
\tikzstyle{arrow} = [thick,->,>=stealth]
\ \medskip

\centering
\begin{tikzpicture}[node distance=35mm and 100mm]
	\node (pu) [startstop] {Prima Única del servicio\\ anotado \(PU(.)\) o \(PV_0(.)\)};
	\node  [below left of=pu, text width=55mm, node distance=18mm, xshift=12mm] {El punto es \(A\) o \(a\) o \(E\) o ...\\
	avec \(x,k,n,m,K...\)};
	\node (puu) [startstop, right of=pu, node distance=50mm] {Valor de una prima periódica unitaria\\ anotada \(PPU_0\) o \(VP_P\)};
	\node  [below of=puu, text width=3cm, node distance=12mm] { \(PPU=_{m'|}\ddot{a}_{x:\lcroof{n'}}^{(k')} \)};
	\node (ppp) [startstop, below of=pu] {Prime Pure Périodique\\ \(PPP\)};
	\node  [below right of=ppp, text width=3cm, node distance=25mm, yshift=6mm] {\(\displaystyle PPP=\frac{PU(.)}{PPU_0}\)};
	\node (ppc) [startstop, right of=ppp, node distance=50mm] {Prima Periódica Cobrada \(PPC\)};
	\node  [below of=ppc, node distance=12mm] {\(\displaystyle PPC=\frac{PPP(1+\lambda)}{1-TxCom}\)};
	\node (Prov) [startstop, below of=ppp] {Reservas matemáticas \(t^{th}\) año\\ indicadas como \(PM_t\) o \(V_t\)};
	\node (provt) [below right of=Prov,text width=10cm, node distance=30mm] {\(V_t=PV_t(.)-PPP\times PPU_t\) \\
	\(PV_t(.)\) recalculado con \(x+t,k,n,m-t,K...\ si\ t\geq m\) \\
	o \(x+t,k,n-(t-m),m=0,K...\ si\ t>m\)\\
	asimismo \(PPU_t=_{(m'-t)^+|}\ddot{a}_{x+t:\lcroof{n'-(t-m)^+}}^{(k')} \)};
%	\node (pro2b) [startstop, right of=dec1, xshift=2cm] {Process 2b};
%	\node (stop) [startstop, below of=pro2b] {Stop};
	
	\draw [arrow] (pu) -- (ppp);
	\draw [arrow] (puu) -- (ppp);
	\draw [arrow] (ppp) -- (ppc);
	\draw [arrow] (ppp) -- node[anchor=east] {\(t\)} (Prov);
%	\draw [arrow] (dec1) -- node[anchor=south] {no} (pro2b);
%	\draw [arrow] (pro2b) |- (ppp);
%	\draw [arrow] (ppp) -- (stop);
	
\end{tikzpicture}
\end{f}